\documentclass[11pt]{article}
\usepackage[margin=1in]{geometry}
\usepackage{amsmath,amssymb}
\usepackage[colorlinks=true,linkcolor=blue,urlcolor=blue,citecolor=blue]{hyperref}
\usepackage{microtype}

\title{The nonabelian Fourier--Stieltjes idempotent norm gap}
\author{Literature-implied answer note for arXiv:math/0405063}
\date{11 August 2026}

\begin{document}
\maketitle

\noindent\textbf{Status.}
\texttt{literature\_implied\_answer (full question)}.

\section*{Source question}
Ilie and Spronk \cite{IlieSpronk} prove that a nonzero idempotent
$u\in B(G)$ has norm one exactly when it is the indicator of an open coset.
Immediately after their Theorem~2.1, on journal page~488 (PDF page~9), they
recall Saeki's abelian gap
\[
  \|u\|_{B(G)}=1
  \quad\text{or}\quad
  \|u\|_{B(G)}\geq c_0,
  \qquad c_0=\frac{1+\sqrt2}{2},
\]
and write: ``We do not know if a similar result holds for nonabelian groups.''

\section*{Supporting theorem}
Mudge and Pham \cite[Theorem~2.2]{MudgePham} prove a stronger statement.
For every locally compact group $G$, if a nonzero idempotent
$u=\chi_S\in M_{cb}A(G)$ satisfies
\[
  \|u\|_{cb}<c_0,
\]
then $S$ is an open coset.  The theorem is on PDF page~3.  Their introduction,
on PDF page~2, records the contractive inclusion
\[
  B(G)\subseteq M_{cb}A(G),
  \qquad \|v\|_{cb}\leq \|v\|_{B(G)} \quad (v\in B(G)).
\]

\section*{Identification}
Let $u\in B(G)$ be a nonzero idempotent with $\|u\|_{B(G)}<c_0$.
The norm comparison gives $\|u\|_{cb}<c_0$, so Mudge--Pham's theorem makes
$u$ the indicator of an open coset.  Ilie--Spronk's Theorem~2.1(i) then gives
$\|u\|_{B(G)}=1$.  Therefore, for every locally compact group, including every
nonabelian one,
\[
  \boxed{\ \|u\|_{B(G)}=1\quad\text{or}\quad
  \|u\|_{B(G)}\geq\frac{1+\sqrt2}{2}.\ }
\]
The constant is sharp for the class of all locally compact groups because the
abelian examples establishing sharpness remain members of that class.

\section*{Provenance and scope}
This is an agent-identified implication, rather than a new mathematical
result.  The later theorem is stated for the larger algebra $M_{cb}A(G)$ and
therefore answers the source sentence after applying its explicit norm
comparison.  A bounded search used the exact source sentence, the threshold,
the terms ``idempotent norm gap'' and ``nonabelian Fourier--Stieltjes'', and
later citing papers; no stronger interpretation of the source question was
found.  The original question, as stated, has no remaining case.

\begin{thebibliography}{9}
\bibitem{IlieSpronk}
M.~Ilie and N.~Spronk,
\emph{Completely bounded homomorphisms of the Fourier algebras},
J. Funct. Anal. \textbf{225} (2005), 480--499;
\href{https://arxiv.org/abs/math/0405063}{arXiv:math/0405063}.

\bibitem{MudgePham}
J.~Mudge and H.~L.~Pham,
\emph{Idempotents with small norms},
J. Funct. Anal. \textbf{270} (2016), 4597--4603;
\href{https://arxiv.org/abs/1510.03535}{arXiv:1510.03535}.
\end{thebibliography}

\end{document}

